\documentclass[11pt,letterpaper]{article}

\usepackage[utf8]{inputenc}
\usepackage[spanish]{babel}
\usepackage[margin=2.5cm]{geometry}
\usepackage{amsmath,amssymb}
\usepackage{enumitem}
\usepackage{titlesec}
\usepackage{array}
\usepackage{hyperref}
\usepackage{xcolor}

\definecolor{unicaucaverde}{RGB}{0,90,60}
\hypersetup{colorlinks=true,linkcolor=unicaucaverde,urlcolor=unicaucaverde}

\titleformat{\section}{\normalfont\Large\bfseries\color{unicaucaverde}}{\thesection}{1em}{}
\titleformat{\subsection}{\normalfont\large\bfseries}{\thesubsection}{1em}{}

\setlength{\parindent}{0pt}
\setlength{\parskip}{6pt}

\title{\textcolor{unicaucaverde}{\Large Guía de Aprendizaje --- Semana 2}\\[4pt]
  \large Estructuras de control: condicionales y ciclos}
\author{Programación --- Universidad del Cauca}
\date{2026}

\begin{document}
\maketitle

\begin{center}
\begin{tabular}{|>{\bfseries}p{4cm}|p{10cm}|}
\hline
Semana & 2 de 16 \\ \hline
Unidad & Fundamentos de programación en C \\ \hline
Subtemas & Condicionales (\texttt{if}/\texttt{else}/\texttt{switch}), ciclos
  (\texttt{while}/\texttt{do-while}/\texttt{for}), \texttt{break}/\texttt{continue} \\ \hline
Horas de clase & 4 \\ \hline
Resultado de aprendizaje & RA1 (parte 2 de 2; continúa de Semana 1) \\ \hline
\end{tabular}
\end{center}

\section*{Relevancia para ingeniería}

Casi ningún programa útil ejecuta sus instrucciones en línea recta de principio a fin: necesita
tomar decisiones (¿el sensor superó el umbral?) y repetir tareas (leer 100 muestras, iterar hasta
converger). Las estructuras de control --- condicionales y ciclos --- son las piezas que permiten
que un programa reaccione a los datos en vez de solo transformarlos. Todo algoritmo de control,
procesamiento de señales o análisis numérico que verá el estudiante más adelante se construye sobre
estas dos ideas.

\section{Objetivo de la semana}

Diseñar y escribir programas en C que utilicen estructuras condicionales y estructuras cíclicas
para resolver problemas que requieren toma de decisiones y repetición de instrucciones.

\section{Resultados de aprendizaje de la semana}

\begin{itemize}[leftmargin=1.2cm]
  \item[\textbf{RA2.1.}] Usar \texttt{if}, \texttt{if-else}, \texttt{else if} y \texttt{switch-case}
        para tomar decisiones en un programa, incluyendo condiciones compuestas con operadores
        lógicos.
  \item[\textbf{RA2.2.}] Usar \texttt{while}, \texttt{do-while} y \texttt{for} --- incluyendo
        ciclos anidados y las instrucciones \texttt{break}/\texttt{continue} --- para repetir
        bloques de instrucciones, eligiendo la estructura adecuada según el problema.
\end{itemize}

\section{Contenidos}

\begin{itemize}[leftmargin=1.2cm]
  \item \textbf{2.1.} Estructuras condicionales: \texttt{if}, \texttt{if-else}, \texttt{else if},
        operador ternario, \texttt{switch-case}.
  \item \textbf{2.2.} Ciclos: \texttt{while}, \texttt{do-while}, \texttt{for};
        \texttt{break}/\texttt{continue}; anidamiento.
\end{itemize}

\section{Plan de la sesión (4 horas)}

\begin{enumerate}[leftmargin=1.2cm]
  \item[\textbf{Hora 1.}] \texttt{if}, \texttt{if-else}, \texttt{else if} y operador ternario;
        ejemplos (clasificar un número, aprobado/reprobado según una nota) (teoría + ejemplos).
  \item[\textbf{Hora 2.}] Taller: condiciones compuestas con \texttt{\&\& || !}; construcción de un
        menú de opciones con \texttt{switch-case} (por ejemplo, una calculadora básica).
  \item[\textbf{Hora 3.}] \texttt{while}, \texttt{do-while} y \texttt{for}: diferencias, cuándo usar
        cada uno, y las instrucciones \texttt{break}/\texttt{continue} (teoría + ejemplos).
  \item[\textbf{Hora 4.}] Taller: ciclos anidados (tabla de multiplicar, patrones con asteriscos),
        contadores y acumuladores, validación de entrada de usuario con \texttt{while}.
\end{enumerate}

\section{Actividades}

\begin{itemize}[leftmargin=1.2cm]
  \item \textbf{En clase}: talleres de las horas 2 y 4 (ver arriba).
  \item \textbf{Trabajo autónomo}: ejercicios propuestos de \texttt{notas.tex} §5 y refuerzo de
        \texttt{ejercicios.tex}.
  \item \textbf{Software}: un compilador de C (\texttt{gcc} local o compilador en línea) para
        escribir, compilar y ejecutar cada ejemplo y ejercicio.
\end{itemize}

\section{Material de apoyo}

\begin{itemize}[leftmargin=1.2cm]
  \item \texttt{notas.tex} --- desarrollo teórico completo de 2.1 y 2.2, con ejemplos de código y
        ejercicios resueltos.
  \item \texttt{diapositivas.tex} --- síntesis visual de la sesión.
  \item \texttt{ejercicios.tex} --- lista adicional de ejercicios de refuerzo con soluciones
        completas en código C.
\end{itemize}

\section{Evaluación de la semana}

Control formativo corto al inicio de la Semana 3 (10 minutos): trazar a mano la salida de un
programa corto en C con un condicional y un ciclo, para verificar RA2.1 y RA2.2 antes de avanzar a
la siguiente unidad.

\section{Bibliografía de la semana}

\begin{enumerate}[leftmargin=1.2cm]
  \item Deitel, P., Deitel, H. \emph{C: How to Program}. 8a ed., Pearson, cap. 3--4.
  \item Kernighan, B. W., Ritchie, D. M. \emph{The C Programming Language}. 2a ed., Prentice Hall,
        cap. 3.
  \item Cairo, L. J. \emph{Fundamentos de programación}. 4a ed., McGraw-Hill.
\end{enumerate}

\end{document}
